\chapter{Communication Complexity — Private Coin Composition}
%%% generated by the blueprint pipeline from TCSlib.CommunicationComplexity.NewmanTheorem.PrivateCoinComposition
\section{Overview}

This module develops composition operations for private-coin communication protocols
with finite message sets.  Starting from their deterministic counterparts, it lifts
\texttt{CommunicationComplexity.Deterministic.FiniteMessage.Protocol.map}, \texttt{CommunicationComplexity.Deterministic.FiniteMessage.Protocol.bind}, \texttt{CommunicationComplexity.PrivateCoin.FiniteMessage.Protocol.comapRandomness}, \texttt{CommunicationComplexity.Deterministic.FiniteMessage.Protocol.prod}, \texttt{CommunicationComplexity.Deterministic.FiniteMessage.Protocol.pi}, \texttt{CommunicationComplexity.PrivateCoin.FiniteMessage.Protocol.rbind},
and \texttt{prodDet} to the private-coin setting, together with the corresponding
simulation lemmas (\texttt{rrun}) and complexity-preservation theorems.

%%%%%%%%%%%%%%%%%%%%%%%%%%%%%%%%%%%%%%%%%%%%%%%%%%%%%%%%%%%%%%%%%%%%%%%%%%%%%%%
\section{Declarations}

\begin{definition}[Output map of a private-coin protocol]
\lean{CommunicationComplexity.PrivateCoin.FiniteMessage.Protocol.map}
\label{CommunicationComplexity.PrivateCoin.FiniteMessage.Protocol.map}
\leanok
\uses{CommunicationComplexity.Deterministic.FiniteMessage.Protocol.map, CommunicationComplexity.PrivateCoin.FiniteMessage.Protocol}
Given a function $g : \alpha \to \beta$ and a private-coin protocol $p$ over randomness
spaces $\Omega_X, \Omega_Y$ with output type $\alpha$, \texttt{map}~$g$~$p$ is the
protocol that runs $p$ and then applies $g$ to its output, yielding a protocol with
output type $\beta$ and the same communication tree.
\end{definition}

\begin{theorem}[Mapping commutes with private-coin execution]
\lean{CommunicationComplexity.PrivateCoin.FiniteMessage.Protocol.map_rrun}
\label{CommunicationComplexity.PrivateCoin.FiniteMessage.Protocol.map_rrun}
\leanok
\uses{CommunicationComplexity.Deterministic.FiniteMessage.Protocol.map_run, CommunicationComplexity.PrivateCoin.FiniteMessage.Protocol, CommunicationComplexity.PrivateCoin.FiniteMessage.Protocol.map, CommunicationComplexity.PrivateCoin.FiniteMessage.Protocol.rrun}
\difficulty{1}
Let $p$ be a private-coin finite-message communication protocol with input types $X$ and
$Y$, private-randomness types $\Omega_X$ and $\Omega_Y$, and output type $\alpha$, and
let $g \colon \alpha \to \beta$. Write $\mathrm{map}(g, p)$ for the protocol that runs
$p$ and then applies $g$ to the result, and let executing a protocol on inputs $x, y$
with private coins $\omega_x, \omega_y$ denote running it on the paired inputs
$(\omega_x, x)$ and $(\omega_y, y)$. Then for all $x \in X$, $y \in Y$, $\omega_x \in
\Omega_X$, and $\omega_y \in \Omega_Y$,
\[
\bigl(\text{executing } \mathrm{map}(g, p) \text{ on } x, y \text{ with coins }
\omega_x, \omega_y\bigr)
  \;=\;
g\bigl(\text{executing } p \text{ on } x, y \text{ with coins } \omega_x,
\omega_y\bigr).
\]
\end{theorem}

\begin{theorem}[Output map preserves communication complexity]
\lean{CommunicationComplexity.PrivateCoin.FiniteMessage.Protocol.map_complexity}
\label{CommunicationComplexity.PrivateCoin.FiniteMessage.Protocol.map_complexity}
\leanok
\uses{CommunicationComplexity.Deterministic.FiniteMessage.Protocol.complexity, CommunicationComplexity.Deterministic.FiniteMessage.Protocol.map_complexity, CommunicationComplexity.PrivateCoin.FiniteMessage.Protocol, CommunicationComplexity.PrivateCoin.FiniteMessage.Protocol.map}
\difficulty{1}
Let $p$ be a private-coin finite-message protocol with input types $X$, $Y$,
private-randomness types $\Omega_X$, $\Omega_Y$, and output type $\alpha$, and let $g
\colon \alpha \to \beta$ be any function. Then the protocol obtained by running $p$ and
applying $g$ to its output has the same communication complexity as $p$, i.e.
$\mathrm{complexity}(\mathrm{map}\,g\,p) = \mathrm{complexity}(p)$.
\end{theorem}

\begin{definition}[Monadic bind of private-coin protocols]
\lean{CommunicationComplexity.PrivateCoin.FiniteMessage.Protocol.bind}
\label{CommunicationComplexity.PrivateCoin.FiniteMessage.Protocol.bind}
\leanok
\uses{CommunicationComplexity.Deterministic.FiniteMessage.Protocol.bind, CommunicationComplexity.PrivateCoin.FiniteMessage.Protocol}
Given a protocol $p$ with output type $\alpha$ and a family $q : \alpha \to \mathrm{Protocol}\;\Omega_X\;\Omega_Y\;X\;Y\;\beta$,
\texttt{CommunicationComplexity.Deterministic.FiniteMessage.Protocol.bind}~$p$~$q$ is the protocol that runs $p$, observes its output $a$, and then
executes $q\,a$ on the same inputs using the same randomness spaces.
\end{definition}

\begin{theorem}[Execution of a sequentially composed private-coin protocol]
\lean{CommunicationComplexity.PrivateCoin.FiniteMessage.Protocol.bind_rrun}
\label{CommunicationComplexity.PrivateCoin.FiniteMessage.Protocol.bind_rrun}
\leanok
\uses{CommunicationComplexity.Deterministic.FiniteMessage.Protocol.bind_run, CommunicationComplexity.PrivateCoin.FiniteMessage.Protocol, CommunicationComplexity.PrivateCoin.FiniteMessage.Protocol.bind, CommunicationComplexity.PrivateCoin.FiniteMessage.Protocol.rrun}
\difficulty{1}
Let $p$ be a private-coin finite-message protocol with input types $X$, $Y$,
private-coin types $\Omega_X$, $\Omega_Y$, and output type $\alpha$, and let $q$ assign
to each value $a \in \alpha$ a private-coin finite-message protocol $q(a)$ over input
types $X$, $Y$ and private-coin types $\Omega_X$, $\Omega_Y$ with output type $\beta$.
Then for all inputs $x \in X$, $y \in Y$ and private coins $\omega_x \in \Omega_X$,
$\omega_y \in \Omega_Y$,
\[
  (p \mathbin{\gg\!=} q)\;\text{run on }(x,y)\text{ with coins }(\omega_x,\omega_y)
  \;=\;
  q(a)\;\text{run on }(x,y)\text{ with coins }(\omega_x,\omega_y),
\]
where $a \in \alpha$ is the output of executing $p$ on inputs $x, y$ with coins
$\omega_x, \omega_y$. That is, running the sequential composition of $p$ and $q$ on
inputs $x, y$ with private coins $\omega_x, \omega_y$ gives the value obtained by first
running $p$ on $x, y$ with coins $\omega_x, \omega_y$ to produce an output $a$, and then
running $q(a)$ on $x, y$ with coins $\omega_x, \omega_y$.
\end{theorem}

\begin{definition}[Randomness reindexing]
\lean{CommunicationComplexity.PrivateCoin.FiniteMessage.Protocol.comapRandomness}
\label{CommunicationComplexity.PrivateCoin.FiniteMessage.Protocol.comapRandomness}
\leanok
\uses{CommunicationComplexity.Deterministic.FiniteMessage.Protocol.comap, CommunicationComplexity.PrivateCoin.FiniteMessage.Protocol}
Given functions $hX : \Omega_X' \to \Omega_X$ and $hY : \Omega_Y' \to \Omega_Y$,
\texttt{CommunicationComplexity.PrivateCoin.FiniteMessage.Protocol.comapRandomness}~$hX$~$hY$~$p$ reindexes the randomness spaces of $p$, yielding
a protocol over $\Omega_X'$ and $\Omega_Y'$ that maps coins through $hX$ and $hY$
before passing them to $p$.
\end{definition}

\begin{theorem}[Randomness reindexing commutes with execution]
\lean{CommunicationComplexity.PrivateCoin.FiniteMessage.Protocol.comapRandomness_rrun}
\label{CommunicationComplexity.PrivateCoin.FiniteMessage.Protocol.comapRandomness_rrun}
\leanok
\uses{CommunicationComplexity.Deterministic.FiniteMessage.Protocol.comap_run, CommunicationComplexity.PrivateCoin.FiniteMessage.Protocol, CommunicationComplexity.PrivateCoin.FiniteMessage.Protocol.comapRandomness, CommunicationComplexity.PrivateCoin.FiniteMessage.Protocol.rrun}
\difficulty{1}
Let $p$ be a private-coin finite-message protocol over input types $X$, $Y$ with
private-randomness types $\Omega_X$, $\Omega_Y$ and output type $\alpha$, and let $hX :
\Omega_X' \to \Omega_X$ and $hY : \Omega_Y' \to \Omega_Y$ be functions reindexing the
two randomness spaces. Then for all inputs $x : X$, $y : Y$ and all coins $\omega_x :
\Omega_X'$, $\omega_y : \Omega_Y'$, executing the reindexed protocol on $x, y$ with
coins $\omega_x, \omega_y$ yields the same output as executing $p$ itself on $x, y$ with
the pushed-forward coins $hX(\omega_x)$ and $hY(\omega_y)$.
\end{theorem}

\begin{theorem}[Complexity is invariant under randomness reindexing]
\lean{CommunicationComplexity.PrivateCoin.FiniteMessage.Protocol.comapRandomness_complexity}
\label{CommunicationComplexity.PrivateCoin.FiniteMessage.Protocol.comapRandomness_complexity}
\leanok
\uses{CommunicationComplexity.Deterministic.FiniteMessage.Protocol.complexity, CommunicationComplexity.PrivateCoin.FiniteMessage.Protocol, CommunicationComplexity.PrivateCoin.FiniteMessage.Protocol.comapRandomness}
\difficulty{1}
Let $p$ be a private-coin finite-message protocol over input types $X$, $Y$ with
private-randomness types $\Omega_X$, $\Omega_Y$ and output type $\alpha$, and let $hX :
\Omega_X' \to \Omega_X$ and $hY : \Omega_Y' \to \Omega_Y$ be arbitrary maps of
randomness spaces. Reindexing the randomness of $p$ through $hX$ and $hY$ — that is,
forming the protocol over $\Omega_X'$, $\Omega_Y'$ that maps its coins through $hX$ and
$hY$ before feeding them to $p$ — yields a protocol with the same worst-case
communication complexity as $p$.
\end{theorem}

\begin{definition}[Product of two private-coin protocols]
\lean{CommunicationComplexity.PrivateCoin.FiniteMessage.Protocol.prod}
\label{CommunicationComplexity.PrivateCoin.FiniteMessage.Protocol.prod}
\leanok
\uses{CommunicationComplexity.Deterministic.FiniteMessage.Protocol.comap, CommunicationComplexity.Deterministic.FiniteMessage.Protocol.prod, CommunicationComplexity.PrivateCoin.FiniteMessage.Protocol}
Given protocols $p_1$ and $p_2$ with potentially distinct randomness spaces
$(\Omega_{X_1}, \Omega_{Y_1})$ and $(\Omega_{X_2}, \Omega_{Y_2})$,
\texttt{CommunicationComplexity.Deterministic.FiniteMessage.Protocol.prod}~$p_1$~$p_2$ runs them in parallel on the same inputs using independent
randomness drawn from $\Omega_{X_1} \times \Omega_{X_2}$ and $\Omega_{Y_1} \times \Omega_{Y_2}$,
returning the pair of their outputs in $\alpha_1 \times \alpha_2$.
\end{definition}

\begin{theorem}[Product protocol runs its factors in parallel]
\lean{CommunicationComplexity.PrivateCoin.FiniteMessage.Protocol.prod_rrun}
\label{CommunicationComplexity.PrivateCoin.FiniteMessage.Protocol.prod_rrun}
\leanok
\uses{CommunicationComplexity.Deterministic.FiniteMessage.Protocol.comap_run, CommunicationComplexity.Deterministic.FiniteMessage.Protocol.prod_run, CommunicationComplexity.PrivateCoin.FiniteMessage.Protocol, CommunicationComplexity.PrivateCoin.FiniteMessage.Protocol.prod, CommunicationComplexity.PrivateCoin.FiniteMessage.Protocol.rrun}
\difficulty{1}
Let $p_1$ and $p_2$ be private-coin finite-message protocols over the same input types
$X$ and $Y$, with output types $\alpha_1$ and $\alpha_2$ and private-randomness types
$(\Omega_{X_1}, \Omega_{Y_1})$ and $(\Omega_{X_2}, \Omega_{Y_2})$ respectively. Then for
all inputs $x : X$ and $y : Y$, and all paired private coins $\omega_x = (\omega_{x,1},
\omega_{x,2}) \in \Omega_{X_1} \times \Omega_{X_2}$ and $\omega_y = (\omega_{y,1},
\omega_{y,2}) \in \Omega_{Y_1} \times \Omega_{Y_2}$, executing the product protocol of
$p_1$ and $p_2$ yields the pair
\[
\bigl(p_1(x, y; \omega_{x,1}, \omega_{y,1}),\; p_2(x, y; \omega_{x,2},
\omega_{y,2})\bigr),
\]
where each component is the output of the corresponding factor run on the same inputs
with its own share of the private randomness.
\end{theorem}

\begin{theorem}[Communication cost of the product protocol]
\lean{CommunicationComplexity.PrivateCoin.FiniteMessage.Protocol.prod_complexity}
\label{CommunicationComplexity.PrivateCoin.FiniteMessage.Protocol.prod_complexity}
\leanok
\uses{CommunicationComplexity.Deterministic.FiniteMessage.Protocol.complexity, CommunicationComplexity.Deterministic.FiniteMessage.Protocol.prod_complexity, CommunicationComplexity.PrivateCoin.FiniteMessage.Protocol, CommunicationComplexity.PrivateCoin.FiniteMessage.Protocol.prod}
\difficulty{1}
Let $p_1$ and $p_2$ be private-coin finite-message protocols, both over the same input
types $X$ and $Y$, with private-randomness types $(\Omega_{X_1}, \Omega_{Y_1})$ and
$(\Omega_{X_2}, \Omega_{Y_2})$ and output types $\alpha_1$ and $\alpha_2$ respectively.
Their product, the protocol on inputs $X$ and $Y$ that runs $p_1$ and $p_2$ in parallel
using independent randomness drawn from $\Omega_{X_1} \times \Omega_{X_2}$ and
$\Omega_{Y_1} \times \Omega_{Y_2}$ and returns the pair of their outputs, has worst-case
communication cost (in bits) equal to the sum of the individual costs:
\[
\mathrm{complexity}(p_1 \times p_2) = \mathrm{complexity}(p_1) +
\mathrm{complexity}(p_2).
\]
\end{theorem}

\begin{definition}[$k$-fold product of private-coin protocols]
\lean{CommunicationComplexity.PrivateCoin.FiniteMessage.Protocol.pi}
\label{CommunicationComplexity.PrivateCoin.FiniteMessage.Protocol.pi}
\leanok
\uses{CommunicationComplexity.Deterministic.FiniteMessage.Protocol.comap, CommunicationComplexity.Deterministic.FiniteMessage.Protocol.pi, CommunicationComplexity.PrivateCoin.FiniteMessage.Protocol}
Given a family of protocols $p : (i : \mathrm{Fin}\,k) \to \mathrm{Protocol}\;(\Omega_{Xf}\,i)\;(\Omega_{Yf}\,i)\;X\;Y\;(\alpha_f\,i)$
with heterogeneous randomness and output types, \texttt{pi}~$p$ runs all $k$ protocols
in parallel on the same inputs using independent randomness from the product spaces
$(\prod_i \Omega_{Xf}\,i,\, \prod_i \Omega_{Yf}\,i)$, returning the tuple of outputs
$(i : \mathrm{Fin}\,k) \to \alpha_f\,i$.
\end{definition}

\begin{theorem}[Componentwise execution of the product protocol]
\lean{CommunicationComplexity.PrivateCoin.FiniteMessage.Protocol.pi_rrun}
\label{CommunicationComplexity.PrivateCoin.FiniteMessage.Protocol.pi_rrun}
\leanok
\uses{CommunicationComplexity.Deterministic.FiniteMessage.Protocol.comap_run, CommunicationComplexity.Deterministic.FiniteMessage.Protocol.pi_run, CommunicationComplexity.PrivateCoin.FiniteMessage.Protocol, CommunicationComplexity.PrivateCoin.FiniteMessage.Protocol.pi, CommunicationComplexity.PrivateCoin.FiniteMessage.Protocol.rrun}
\difficulty{1}
Let $k$ be a natural number and, for each $i \in \mathrm{Fin}\,k$, let $p_i$ be a
private-coin finite-message protocol with private-randomness types $\Omega_{X,i}$ for
Alice and $\Omega_{Y,i}$ for Bob, common public input types $X$ and $Y$, and output type
$\alpha_i$. Fix public inputs $x \in X$ and $y \in Y$, together with coin vectors
$\omega_x \in \prod_{i} \Omega_{X,i}$ and $\omega_y \in \prod_{i} \Omega_{Y,i}$. Then
executing the product protocol on these data agrees componentwise with executing each
factor on its own coordinate of the randomness:
\[
  (\pi\,p).\mathrm{rrun}\;x\;y\;\omega_x\;\omega_y
\;=\; \bigl(\, i \mapsto (p_i).\mathrm{rrun}\;x\;y\;(\omega_x\,i)\;(\omega_y\,i)
\,\bigr).
\]
\end{theorem}

\begin{theorem}[Additivity of complexity under products]
\lean{CommunicationComplexity.PrivateCoin.FiniteMessage.Protocol.pi_complexity}
\label{CommunicationComplexity.PrivateCoin.FiniteMessage.Protocol.pi_complexity}
\leanok
\uses{CommunicationComplexity.Deterministic.FiniteMessage.Protocol.complexity, CommunicationComplexity.Deterministic.FiniteMessage.Protocol.pi_complexity, CommunicationComplexity.PrivateCoin.FiniteMessage.Protocol, CommunicationComplexity.PrivateCoin.FiniteMessage.Protocol.pi}
\difficulty{1}
Let $k$ be a natural number and let $p_0, \dots, p_{k-1}$ be a family of private-coin
finite-message protocols indexed by $i \in \mathrm{Fin}\,k$, where $p_i$ has
private-randomness types $\Omega_{X,i}$, $\Omega_{Y,i}$, shared input types $X$, $Y$,
and output type $\alpha_i$. Then the communication complexity of their $k$-fold product
protocol equals the sum of the individual communication complexities, \[
\mathrm{complexity}\bigl(\textstyle\prod_{i} p_i\bigr) = \sum_{i}
\mathrm{complexity}(p_i). \]
\end{theorem}

\begin{definition}[Bind with fresh independent randomness]
\lean{CommunicationComplexity.PrivateCoin.FiniteMessage.Protocol.rbind}
\label{CommunicationComplexity.PrivateCoin.FiniteMessage.Protocol.rbind}
\leanok
\uses{CommunicationComplexity.Deterministic.FiniteMessage.Protocol.bind, CommunicationComplexity.Deterministic.FiniteMessage.Protocol.comap, CommunicationComplexity.PrivateCoin.FiniteMessage.Protocol}
Given a protocol $p$ over randomness $(\Omega_X, \Omega_Y)$ and a family
$q : \alpha \to \mathrm{Protocol}\;\Omega_X'\;\Omega_Y'\;X\;Y\;\beta$,
\texttt{CommunicationComplexity.PrivateCoin.FiniteMessage.Protocol.rbind}~$p$~$q$ runs $p$ with the first component of the coin pair,
obtains its output $a$, then runs $q\,a$ using a fresh, independent second component,
yielding a protocol over $(\Omega_X \times \Omega_X', \Omega_Y \times \Omega_Y')$.
\end{definition}

\begin{theorem}[Fresh-randomness bind computes by splitting the coins]
\lean{CommunicationComplexity.PrivateCoin.FiniteMessage.Protocol.rbind_rrun}
\label{CommunicationComplexity.PrivateCoin.FiniteMessage.Protocol.rbind_rrun}
\leanok
\uses{CommunicationComplexity.Deterministic.FiniteMessage.Protocol.bind_run, CommunicationComplexity.Deterministic.FiniteMessage.Protocol.comap_run, CommunicationComplexity.PrivateCoin.FiniteMessage.Protocol, CommunicationComplexity.PrivateCoin.FiniteMessage.Protocol.rbind, CommunicationComplexity.PrivateCoin.FiniteMessage.Protocol.rrun}
\difficulty{1}
Let $p$ be a private-coin finite-message protocol over inputs $X$, $Y$ with private
randomness types $\Omega_X$, $\Omega_Y$ and output type $\alpha$, and let $q$ assign to
each output $a\in\alpha$ a private-coin finite-message protocol $q(a)$ over the same
inputs $X$, $Y$ with fresh randomness types $\Omega_X'$, $\Omega_Y'$ and output type
$\beta$. Then for all inputs $x\in X$, $y\in Y$ and all coin pairs
$\omega_x=(\omega_x^1,\omega_x^2)\in\Omega_X\times\Omega_X'$ and
$\omega_y=(\omega_y^1,\omega_y^2)\in\Omega_Y\times\Omega_Y'$, executing the
fresh-randomness bind of $p$ and $q$ agrees with running $p$ on the first coordinates of
the coins to obtain an output $a$ and then running $q(a)$ on the second coordinates:
\[
  \mathrm{rrun}\big(\mathrm{rbind}(p,q)\big)\,x\,y\,\omega_x\,\omega_y
  \;=\;
\mathrm{rrun}\big(q(\mathrm{rrun}\,p\,x\,y\,\omega_x^1\,\omega_y^1)\big)\,x\,y\,\omega_x^2\,\omega_y^2
.
\]
\end{theorem}

\begin{theorem}[Complexity of bind with fresh randomness under a constant continuation]
\lean{CommunicationComplexity.PrivateCoin.FiniteMessage.Protocol.rbind_complexity_const}
\label{CommunicationComplexity.PrivateCoin.FiniteMessage.Protocol.rbind_complexity_const}
\leanok
\uses{CommunicationComplexity.Deterministic.FiniteMessage.Protocol.bind_complexity_const, CommunicationComplexity.Deterministic.FiniteMessage.Protocol.complexity, CommunicationComplexity.PrivateCoin.FiniteMessage.Protocol, CommunicationComplexity.PrivateCoin.FiniteMessage.Protocol.rbind}
\difficulty{2}
Let $p$ be a private-coin finite-message protocol over inputs $X$, $Y$ with private
randomness $(\Omega_X, \Omega_Y)$ and output type $\alpha$, and let $q$ assign to each
output $a \in \alpha$ a private-coin finite-message protocol $q\,a$ over $X$, $Y$ with
fresh randomness $(\Omega_X', \Omega_Y')$ and output type $\beta$. Suppose there is a
natural number $c$ with $\mathrm{complexity}(q\,a) = c$ for every $a \in \alpha$. Then
the protocol obtained by running $p$ and then $q\,a$ with independent randomness
satisfies
\[
  \mathrm{complexity}(\mathrm{rbind}\,p\,q) \;=\; \mathrm{complexity}(p) + c.
\]
\end{theorem}

\begin{definition}[Product of private-coin and deterministic protocols]
\lean{CommunicationComplexity.PrivateCoin.FiniteMessage.Protocol.prodDet}
\label{CommunicationComplexity.PrivateCoin.FiniteMessage.Protocol.prodDet}
\leanok
\uses{CommunicationComplexity.Deterministic.FiniteMessage.Protocol, CommunicationComplexity.Deterministic.FiniteMessage.Protocol.comap, CommunicationComplexity.Deterministic.FiniteMessage.Protocol.map, CommunicationComplexity.PrivateCoin.FiniteMessage.Protocol, CommunicationComplexity.PrivateCoin.FiniteMessage.Protocol.bind}
Given a private-coin protocol $p_1$ and a deterministic protocol $p_2$,
\texttt{CommunicationComplexity.PrivateCoin.FiniteMessage.Protocol.prodDet}~$p_1$~$p_2$ runs both on the same inputs and pairs their outputs
in $\alpha_1 \times \alpha_2$, with $p_2$ contributing no additional randomness.
\end{definition}

\begin{theorem}[Execution of a private-coin–deterministic product protocol]
\lean{CommunicationComplexity.PrivateCoin.FiniteMessage.Protocol.prodDet_rrun}
\label{CommunicationComplexity.PrivateCoin.FiniteMessage.Protocol.prodDet_rrun}
\leanok
\uses{CommunicationComplexity.Deterministic.FiniteMessage.Protocol, CommunicationComplexity.Deterministic.FiniteMessage.Protocol.bind_run, CommunicationComplexity.Deterministic.FiniteMessage.Protocol.comap_run, CommunicationComplexity.Deterministic.FiniteMessage.Protocol.map_run, CommunicationComplexity.Deterministic.FiniteMessage.Protocol.run, CommunicationComplexity.PrivateCoin.FiniteMessage.Protocol, CommunicationComplexity.PrivateCoin.FiniteMessage.Protocol.prodDet, CommunicationComplexity.PrivateCoin.FiniteMessage.Protocol.rrun}
\difficulty{1}
Let $X$, $Y$ be input types and $\Omega_X$, $\Omega_Y$ private-randomness types. Let
$p_1$ be a private-coin finite-message protocol over these types with output type
$\alpha_1$, and let $p_2$ be a deterministic finite-message protocol over $X$, $Y$ with
output type $\alpha_2$. Then for all inputs $x : X$, $y : Y$ and all private coins
$\omega_x : \Omega_X$, $\omega_y : \Omega_Y$, executing the product protocol of $p_1$
and $p_2$ on $x$, $y$ with coins $\omega_x$, $\omega_y$ yields the pair whose first
component is the output of $p_1$ executed on $x$, $y$ with coins $\omega_x$, $\omega_y$
and whose second component is the output of the deterministic protocol $p_2$ run on $x$,
$y$.
\end{theorem}

\begin{theorem}[Additivity of complexity under the product of protocols]
\lean{CommunicationComplexity.PrivateCoin.FiniteMessage.Protocol.prodDet_complexity}
\label{CommunicationComplexity.PrivateCoin.FiniteMessage.Protocol.prodDet_complexity}
\leanok
\uses{CommunicationComplexity.Deterministic.FiniteMessage.Protocol, CommunicationComplexity.Deterministic.FiniteMessage.Protocol.bind_complexity_const, CommunicationComplexity.Deterministic.FiniteMessage.Protocol.complexity, CommunicationComplexity.PrivateCoin.FiniteMessage.Protocol, CommunicationComplexity.PrivateCoin.FiniteMessage.Protocol.prodDet}
\difficulty{2}
Fix input types $X$, $Y$, a private-randomness type $\Omega_X$ for Alice and $\Omega_Y$
for Bob, and output types $\alpha_1$, $\alpha_2$. Let $p_1$ be a private-coin
finite-message protocol producing an output in $\alpha_1$, and let $p_2$ be a
deterministic finite-message protocol over $X$, $Y$ producing an output in $\alpha_2$.
Form their product, the private-coin protocol that runs $p_1$ and then $p_2$ on the same
inputs (with $p_2$ using no private randomness) and reports the pair of their outputs in
$\alpha_1 \times \alpha_2$. Then the worst-case communication cost of this product
equals the sum of the individual costs, that is, the complexity of the product is the
complexity of $p_1$ plus the complexity of $p_2$.
\end{theorem}
